% !PRINTED 544
\weberpage{544}

Macht man auf der rechten Seite dieser Formel in den Indices der $\alpha$ die Permutation $\pi$, so ergiebt sich nach 1.
\[
\sum^{\varepsilon} \varepsilon^{-k-\nu}(\varepsilon,\alpha)^{\nu},
\]
d. h. $A_k^{(\nu)}$ geht in $A_{k+\nu}^{(\nu)}$ über, wie im Satz 2. behauptet ist.

Der Satz 2. ist ein specieller Fall eines allgemeinen Theorems. Entwickeln wir ein Product von beliebig vielen Factoren
\[
(\varepsilon,\alpha)^{\nu}(\varepsilon^{\lambda_1},\alpha)^{\nu_1}(\varepsilon^{\lambda_2},\alpha)^{\nu_2}\dots,
\]
worin $\nu, \nu_1, \nu_2\dots$ positive, $\lambda_1, \lambda_2\dots$ beliebige ganze Zahlen sind, und ordnen es, wie vorher $(\varepsilon,\alpha)^{\nu}$, nach Potenzen von $\varepsilon$, so mag sich ergeben
\weq{7}{(\varepsilon,\alpha)^{\nu}(\varepsilon^{\lambda_1},\alpha)^{\nu_1}(\varepsilon^{\lambda_2},\alpha)^{\nu_2}\dots = \sum_{0,\,m-1}^{h}\varepsilon^h B_h,}

worin die $B_h$ von $\varepsilon$ unabhängige Formen von der Variablen $\alpha$ sind. Da auch diese Entwickelung für alle $m^{\mathrm{ten}}$ Einheitswurzeln $\varepsilon$ gilt, so kann man die Formel (2) anwenden und erhält
% [sic] authority prints B_k rather than mB_k and puts each \lambda_j on \alpha in the
% factors below; coefficient extraction from (7) and the product displayed above require
% mB_k and (\varepsilon^{\lambda_j},\alpha), respectively.
\[
B_k = \sum^{\varepsilon} \varepsilon^{-k}(\varepsilon,\alpha)^{\nu}(\varepsilon,\alpha^{\lambda_1})^{\nu_1}(\varepsilon,\alpha^{\lambda_2})^{\nu_2}\dots
\]
woraus man nach 1. schliessen kann, dass $B_k$ durch die Substitution $\pi$ in
\[
B_{k+\nu+\lambda_1\nu_1+\lambda_2\nu_2+\cdots}
\]
übergeht. Wir haben also:

\gesperrt{3. Die Permutation $\pi$, auf die Indices der $\alpha$ angewandt, ruft unter den Indices der Coëfficienten des nach $\varepsilon$ geordneten Productes}
\weq{8}{(\varepsilon,\alpha)^{\nu}(\varepsilon^{\lambda_1},\alpha)^{\nu_1}(\varepsilon^{2},\alpha)^{\nu_2}\dots} % [sic] authority prints \varepsilon^2 here, omitting the expected \lambda_2.
\gesperrt{die Permutation $\pi^{\nu+\lambda_1\nu_1+\lambda_2\nu_2\cdots}$ hervor.}

In diesen Theoremen kann die Permutation $\pi$ wiederholt werden; sie bleiben richtig, wenn $\pi$ durch irgend eine Potenz von $\pi$ ersetzt wird und sie gelten, was auch $\varepsilon$ für eine $m^{\mathrm{te}}$ Einheitswurzel sein mag. Ist aber $\varepsilon$ eine imprimitive Einheitswurzel, so treten gewisse Vereinfachungen ein, die wir noch kennen lernen müssen. Ist $e$ ein Theiler von $m$,
\[
m = ef
\]
und $\varepsilon$ irgend eine $e^{\mathrm{te}}$ Einheitswurzel, so wird der Ausdruck (1)
\begin{align*}
(\varepsilon,\alpha) &= \alpha + \varepsilon\alpha_1 + \cdots + \varepsilon^{e-1}\alpha_{e-1}\\
&\quad + \alpha_e + \varepsilon\alpha_{e+1} + \cdots + \varepsilon^{e-1}\alpha_{2e-1}\\
&\quad \cdots\cdots\cdots\cdots\cdots\cdots\cdots\cdots\cdots\cdots\\
&\quad + \alpha_{e(f-1)} + \varepsilon\alpha_{e(f-1)+1} + \cdots + \varepsilon^{e-1}\alpha_{m-1}
\end{align*}
